\documentclass{article}
\usepackage{amsmath,amssymb}
\usepackage{xurl}
\usepackage[hidelinks]{hyperref}
% Shared commands for notes in docs/paper-gaps/.

\DeclareUnicodeCharacter{2080}{\ifmmode{}_{0}\else\textsubscript{0}\fi}
\DeclareUnicodeCharacter{2081}{\ifmmode{}_{1}\else\textsubscript{1}\fi}
\DeclareUnicodeCharacter{2082}{\ifmmode{}_{2}\else\textsubscript{2}\fi}
\DeclareUnicodeCharacter{2099}{\ifmmode{}_{n}\else\textsubscript{n}\fi}

\newcommand{\Lean}{\textsc{Lean}}
\ExplSyntaxOn
\NewDocumentCommand{\leanid}{m}{
  \ifmmode
    \textsf{\detokenize{#1}}
  \else
    \group_begin:
    \urlstyle{sf}
    \tl_set:Nn \l_tmpa_tl {#1}
    \tl_replace_all:Nnn \l_tmpa_tl {\_} {_}
    \exp_args:NV \path \l_tmpa_tl
    \group_end:
  \fi
}
\ExplSyntaxOff
\providecommand{\tr}{\operatorname{tr}}
\newcommand{\tnleanIssueBase}{https://github.com/LionSR/TNLean/issues/}
\newcommand{\tnleanPrBase}{https://github.com/LionSR/TNLean/pull/}
\newcommand{\tnleanDocBase}{https://sirui-lu.com/TNLean/docs/}
\newcommand{\ghissue}[1]{\href{\tnleanIssueBase#1}{GitHub issue \##1}}
\newcommand{\ghpr}[1]{\href{\tnleanPrBase#1}{PR \##1}}
\newcommand{\doclink}[1]{\href{\tnleanDocBase#1.html}{\path{#1}}}

% Dirac notation and the identity operator, as in blueprint/src/macros/common.tex.
% \providecommand keeps note-local definitions authoritative.
\usepackage{dsfont}
\providecommand{\ket}[1]{|#1\rangle}
\providecommand{\bra}[1]{\langle #1|}
\providecommand{\braket}[2]{\langle #1 | #2 \rangle}
\providecommand{\ketbra}[2]{|#1\rangle\!\langle #2|}
\providecommand{\Id}{\mathds{1}}
\providecommand{\one}{\mathds{1}}
\providecommand{\C}{\mathbb{C}}
\providecommand{\MN}[1]{M_{#1}(\C)}
\providecommand{\MD}{M_D(\C)}

% Verdict marker: \gapnote{<kind>}{<status>}. Kind is the policy
% classification (clarification, local-correction, scope-restriction,
% unfaithful, false-source, open-gap); status is open, wip (elimination
% underway), resolved, or historical. Machine-read by the site index;
% prints a verdict line.
\newcommand{\gapnote}[2]{\par\noindent\textbf{Verdict:} #1 (#2).\par}


\title{Typed Multiplicity Indices for the Coupled Pentagon}
\author{}
\date{2026-10-05}

\begin{document}
\maketitle
\gapnote{local-correction}{wip}

The action map of Ref.~\cite{GarreRubio2023MPOSym} has type
$V_{ax}^{y,i}:\mathbb C^{\chi_a}\otimes\mathbb C^{D_x}
    \to\mathbb C^{D_y}$, while $\widehat V$ is its synthesis partner.
Consequently the displayed change of action trees is
$H_{\rm seq}=L H_{\rm fus}$ and its coefficient contraction is
$L\otimes I=H_{\rm seq}S_{\rm fus}$. The fusion change uses the opposite
tree direction, $H_{\rm left}=F H_{\rm right}$. These orientations follow
from the typed maps and the diagram legs, without adjoints or transposes.%
\footnote{Source traceability:
    \path{Papers/2203.12563/REsubmission.tex}, action-map type at line 459,
    diagram definitions at lines 97--159, and labels
    \texttt{Fsymbolsdef}, \texttt{eq:F\_symbol2}, and \texttt{1Fsymbol}.}

In the coupled pentagon, the surrounding action factors fix the ranges
\begin{align}
    \eta & \in\{1,\ldots,N_{ab}^{d}\}, &
    \chi & \in\{1,\ldots,N_{dc}^{e}\},   \\
    \mu  & \in\{1,\ldots,N_{bc}^{f}\}, &
    \nu  & \in\{1,\ldots,N_{af}^{e}\}.
    \label{eq:glm_l_index_ranges}
\end{align}
The printed factor
$(F_{abc}^{e})_{f,\nu\mu}^{d,\chi\eta}$ reverses both pairs.
It need not be well typed when these multiplicities differ. The typed
coupled pentagon is
\begin{align}
    \sum_{d,n,\eta,\chi}
    (L_{abt}^{y})_{d,n\eta}^{z,ij}
    (L_{dcx}^{y})_{e,m\chi}^{t,nk}
    (F_{abc}^{e})_{f,\mu\nu}^{d,\eta\chi}
     & ={}
    \sum_l
    (L_{bcx}^{z})_{f,l\mu}^{t,jk}
    (L_{afx}^{y})_{e,m\nu}^{z,il}.
    \label{eq:glm_l_typed_pentagon}
\end{align}
Only the two index orders in the F factor differ from the source's
\texttt{coupledpent}, at lines 554--562.

In the following fusion gauge formula, use fresh labels: $a,b,c$ for the
three input blocks, $d$ for the final block, and $e,f$ for the intermediate
blocks. For analysis-basis gauges $H'_{\rm left}=Y_{\rm left}H_{\rm left}$ and
$H'_{\rm right}=Y_{\rm right}H_{\rm right}$, the fusion law is
\begin{align}
    F' & =Y_{\rm left}F Y_{\rm right}^{-1}, &
    L' & =X_{\rm seq}L X_{\rm fus}^{-1}.
    \label{eq:glm_l_gauge_matrices}
\end{align}
Thus an inverse gauge on the right has its old coordinate as row and its
new coordinate as column. The component display at lines 417--422 also
needs primed indices on the F entry being summed. Its well-typed entry is
\begin{align}
    (F')_{f,\chi\eta}^{e,\mu\nu}
    ={} & \sum_{\mu',\nu',\chi',\eta'}
    (Y_{ab}^{e})_{\mu\mu'}(Y_{ec}^{d})_{\nu\nu'}
    F_{f,\chi'\eta'}^{e,\mu'\nu'}
    (Y_{bc}^{f})^{-1}_{\chi'\chi}
    (Y_{af}^{d})^{-1}_{\eta'\eta}.
    \label{eq:glm_l_typed_f_gauge}
\end{align}
These are ordinary inverse matrices acting on dual coordinate positions,
not adjoints. The source's phrase ``non-zero constants'' means an
invertible change of basis; individual matrix entries may be zero.

The index correction follows from the displayed domains. The multiplicity
L matrices, their inverse identities, and their defining analysis-tree
equation are formalized.\footnote{Formal declarations:
    \leanid{MPOTensor.actionLMatrix_cross},
    \leanid{MPOTensor.actionLMatrix_inverse_of_isInjective}, and
    \leanid{MPOTensor.actionLMatrix_analysis}, in
    \doclink{TNLean/MPS/MPDO/BoundaryActionLMatrix}.}
The corrected coupled pentagon (\ref{eq:glm_l_typed_pentagon}) and the
general multiplicity-gauge law remain separate proof obligations; they are
not assumed by the matrix construction.

\bibliographystyle{alpha}
\bibliography{references}
\end{document}
